% TOP_PROBLEM: {"id": 30006734, "title": "The Borel Conjecture on Topological Rigidity of Aspherical Manifolds", "category_id": 7, "category": {"id": 7, "name": "topology", "display_name": "Topology", "description": "Properties preserved under continuous deformations.", "slug": "topology", "order_index": 7, "created_at": "2026-07-31T15:26:25.671Z"}, "status": "open"}
% FIRST_POSTED: 2026-09-25
% ENTRY_KIND: statement_only
% !TeX program = lualatex
% Definition and sources only; this file is not an LLM solution attempt.
\documentclass[11pt]{article}
\usepackage[margin=25mm]{geometry}
\usepackage{fontspec}
\setmainfont{Latin Modern Roman}
\usepackage{amsmath,amssymb,mathtools}
\usepackage{unicode-math}
\setmathfont{Latin Modern Math}
\usepackage{xurl,hyperref}
\hypersetup{colorlinks=true,urlcolor=blue}
\setcounter{secnumdepth}{0}
\setlength{\parindent}{0pt}
\setlength{\parskip}{5pt}
\setlength{\emergencystretch}{3em}
\begin{document}
\section*{\texorpdfstring{The Borel Conjecture on Topological Rigidity of Aspherical Manifolds}{The Borel Conjecture on Topological Rigidity of Aspherical Manifolds}}\label{30006734}
\subsection{Definitions and mathematical statement}
A connected manifold is aspherical if its universal cover is contractible, equivalently its higher homotopy groups vanish. Given closed aspherical manifolds $M,N$ and a homotopy equivalence $f:M\to N$, the conjecture asks for a homeomorphism $h:M\to N$ with
\[
 f\simeq h.
\]
The specific homotopy class must contain a homeomorphism; merely knowing that the two manifolds happen to be homeomorphic is weaker. The category is topological, so smooth rigidity or a diffeomorphism is not requested.
\par\smallskip\textit{Formulation and concept references:} \href{https://link.springer.com/article/10.1007/s10711-023-00861-4}{[S1]}.

\subsection{Short English statement}
Does every homotopy equivalence between closed aspherical manifolds come from a homeomorphism up to homotopy?
\subsection{Sources}
\begin{itemize}
\item[S1]U. Hamenstadt and F. Jaeckel, Rigidity of geometric structures, Geometriae Dedicata 218, article 16 (2024), Section 2.
\url{https://link.springer.com/article/10.1007/s10711-023-00861-4}.
\textit{Formulation source.}
\end{itemize}
\end{document}
